% Options for packages loaded elsewhere
\PassOptionsToPackage{unicode}{hyperref}
\PassOptionsToPackage{hyphens}{url}
\PassOptionsToPackage{dvipsnames,svgnames,x11names}{xcolor}
\documentclass[
  11pt,
  letterpaper,
]{article}
\usepackage{xcolor}
\usepackage[margin=0.8in]{geometry}
\usepackage{amsmath,amssymb}
\setcounter{secnumdepth}{5}
\usepackage{iftex}
\ifPDFTeX
  \usepackage[T1]{fontenc}
  \usepackage[utf8]{inputenc}
  \usepackage{textcomp} % provide euro and other symbols
\else % if luatex or xetex
  \usepackage{unicode-math} % this also loads fontspec
  \defaultfontfeatures{Scale=MatchLowercase}
  \defaultfontfeatures[\rmfamily]{Ligatures=TeX,Scale=1}
\fi
\usepackage{lmodern}
\ifPDFTeX\else
  % xetex/luatex font selection
\fi
% Use upquote if available, for straight quotes in verbatim environments
\IfFileExists{upquote.sty}{\usepackage{upquote}}{}
\IfFileExists{microtype.sty}{% use microtype if available
  \usepackage[]{microtype}
  \UseMicrotypeSet[protrusion]{basicmath} % disable protrusion for tt fonts
}{}
\makeatletter
\@ifundefined{KOMAClassName}{% if non-KOMA class
  \IfFileExists{parskip.sty}{%
    \usepackage{parskip}
  }{% else
    \setlength{\parindent}{0pt}
    \setlength{\parskip}{6pt plus 2pt minus 1pt}}
}{% if KOMA class
  \KOMAoptions{parskip=half}}
\makeatother
\usepackage{color}
\usepackage{fancyvrb}
\newcommand{\VerbBar}{|}
\newcommand{\VERB}{\Verb[commandchars=\\\{\}]}
\DefineVerbatimEnvironment{Highlighting}{Verbatim}{commandchars=\\\{\}}
% Add ',fontsize=\small' for more characters per line
\newenvironment{Shaded}{}{}
\newcommand{\AlertTok}[1]{\textcolor[rgb]{1.00,0.00,0.00}{\textbf{#1}}}
\newcommand{\AnnotationTok}[1]{\textcolor[rgb]{0.38,0.63,0.69}{\textbf{\textit{#1}}}}
\newcommand{\AttributeTok}[1]{\textcolor[rgb]{0.49,0.56,0.16}{#1}}
\newcommand{\BaseNTok}[1]{\textcolor[rgb]{0.25,0.63,0.44}{#1}}
\newcommand{\BuiltInTok}[1]{\textcolor[rgb]{0.00,0.50,0.00}{#1}}
\newcommand{\CharTok}[1]{\textcolor[rgb]{0.25,0.44,0.63}{#1}}
\newcommand{\CommentTok}[1]{\textcolor[rgb]{0.38,0.63,0.69}{\textit{#1}}}
\newcommand{\CommentVarTok}[1]{\textcolor[rgb]{0.38,0.63,0.69}{\textbf{\textit{#1}}}}
\newcommand{\ConstantTok}[1]{\textcolor[rgb]{0.53,0.00,0.00}{#1}}
\newcommand{\ControlFlowTok}[1]{\textcolor[rgb]{0.00,0.44,0.13}{\textbf{#1}}}
\newcommand{\DataTypeTok}[1]{\textcolor[rgb]{0.56,0.13,0.00}{#1}}
\newcommand{\DecValTok}[1]{\textcolor[rgb]{0.25,0.63,0.44}{#1}}
\newcommand{\DocumentationTok}[1]{\textcolor[rgb]{0.73,0.13,0.13}{\textit{#1}}}
\newcommand{\ErrorTok}[1]{\textcolor[rgb]{1.00,0.00,0.00}{\textbf{#1}}}
\newcommand{\ExtensionTok}[1]{#1}
\newcommand{\FloatTok}[1]{\textcolor[rgb]{0.25,0.63,0.44}{#1}}
\newcommand{\FunctionTok}[1]{\textcolor[rgb]{0.02,0.16,0.49}{#1}}
\newcommand{\ImportTok}[1]{\textcolor[rgb]{0.00,0.50,0.00}{\textbf{#1}}}
\newcommand{\InformationTok}[1]{\textcolor[rgb]{0.38,0.63,0.69}{\textbf{\textit{#1}}}}
\newcommand{\KeywordTok}[1]{\textcolor[rgb]{0.00,0.44,0.13}{\textbf{#1}}}
\newcommand{\NormalTok}[1]{#1}
\newcommand{\OperatorTok}[1]{\textcolor[rgb]{0.40,0.40,0.40}{#1}}
\newcommand{\OtherTok}[1]{\textcolor[rgb]{0.00,0.44,0.13}{#1}}
\newcommand{\PreprocessorTok}[1]{\textcolor[rgb]{0.74,0.48,0.00}{#1}}
\newcommand{\RegionMarkerTok}[1]{#1}
\newcommand{\SpecialCharTok}[1]{\textcolor[rgb]{0.25,0.44,0.63}{#1}}
\newcommand{\SpecialStringTok}[1]{\textcolor[rgb]{0.73,0.40,0.53}{#1}}
\newcommand{\StringTok}[1]{\textcolor[rgb]{0.25,0.44,0.63}{#1}}
\newcommand{\VariableTok}[1]{\textcolor[rgb]{0.10,0.09,0.49}{#1}}
\newcommand{\VerbatimStringTok}[1]{\textcolor[rgb]{0.25,0.44,0.63}{#1}}
\newcommand{\WarningTok}[1]{\textcolor[rgb]{0.38,0.63,0.69}{\textbf{\textit{#1}}}}
\usepackage{longtable,booktabs,array}
\usepackage{calc} % for calculating minipage widths
% Correct order of tables after \paragraph or \subparagraph
\usepackage{etoolbox}
\makeatletter
\patchcmd\longtable{\par}{\if@noskipsec\mbox{}\fi\par}{}{}
\makeatother
% Allow footnotes in longtable head/foot
\IfFileExists{footnotehyper.sty}{\usepackage{footnotehyper}}{\usepackage{footnote}}
\makesavenoteenv{longtable}
\setlength{\emergencystretch}{3em} % prevent overfull lines
\providecommand{\tightlist}{%
  \setlength{\itemsep}{0pt}\setlength{\parskip}{0pt}}
\usepackage{bookmark}
\IfFileExists{xurl.sty}{\usepackage{xurl}}{} % add URL line breaks if available
\urlstyle{same}
\hypersetup{
  pdftitle={Evolutionary Game Modeling of Immigration Policy},
  pdfauthor={Kai Zhang},
  colorlinks=true,
  linkcolor={Maroon},
  filecolor={Maroon},
  citecolor={Blue},
  urlcolor={blue},
  pdfcreator={LaTeX via pandoc}}

\title{Evolutionary Game Modeling of Immigration Policy}
\usepackage{etoolbox}
\makeatletter
\providecommand{\subtitle}[1]{% add subtitle to \maketitle
  \apptocmd{\@title}{\par {\large #1 \par}}{}{}
}
\makeatother
\subtitle{English mathematical notes and reproducibility audit}
\author{Kai Zhang}
\date{23 October 2026}

\begin{document}
\maketitle

{
\setcounter{tocdepth}{3}
\tableofcontents
}
\section{Scope and provenance}\label{scope-and-provenance}

This technical companion reconstructed the mathematical content of three
Python scripts supplied with an undergraduate thesis project. The files
contained \textbf{two different payoff specifications}, not merely
different plotting languages. This document therefore treated them
separately:

\begin{enumerate}
\def\labelenumi{\arabic{enumi}.}
\tightlist
\item
  \textbf{Government-benefit spillovers.}\\
  The scripts \texttt{replicator\_sim\_AB\_no\_pandas.py} and
  \texttt{replicator\_compare\_jp.py} used the same payoffs and a
  clipped fourth-order Runge--Kutta scheme.
\item
  \textbf{Firm-cost spillovers with mutation.}\\
  The script \texttt{replicator\_paper\_ready\_jp.py} used a different
  payoff specification, an effective hiring cost and a small mutation
  term, with clipped explicit Euler updates.
\end{enumerate}

The English exposition, analytic Jacobians, source-table checks and
adaptive-solver diagnostics were prepared for this repository release.
Original scripts and five output tables were preserved under
\texttt{source\_snapshot/} with SHA-256 hashes. No full thesis
manuscript, MATLAB source or empirical calibration dataset was included
in the supplied folder. This document is an English mathematical
reconstruction from the recoverable code, rather than a
sentence-by-sentence translation of an absent manuscript.

\textbf{A consequential numerical finding:} the original \(10^{-9}\)
population floor changed the government model's finite-horizon hiring
path substantially. Reproducing the original tables and solving the
unmodified continuous ODE are different checks; both are reported below.

\section{Players, strategies and state
variables}\label{players-strategies-and-state-variables}

The model represented three interacting populations. A strategy share
was a population-level state, rather than a forecast probability
estimated from data.

\begin{longtable}[]{@{}
  >{\raggedright\arraybackslash}p{(\linewidth - 6\tabcolsep) * \real{0.2500}}
  >{\raggedright\arraybackslash}p{(\linewidth - 6\tabcolsep) * \real{0.2500}}
  >{\raggedright\arraybackslash}p{(\linewidth - 6\tabcolsep) * \real{0.2500}}
  >{\raggedright\arraybackslash}p{(\linewidth - 6\tabcolsep) * \real{0.2500}}@{}}
\toprule\noalign{}
\begin{minipage}[b]{\linewidth}\raggedright
Population
\end{minipage} & \begin{minipage}[b]{\linewidth}\raggedright
First strategy
\end{minipage} & \begin{minipage}[b]{\linewidth}\raggedright
Second strategy
\end{minipage} & \begin{minipage}[b]{\linewidth}\raggedright
State
\end{minipage} \\
\midrule\noalign{}
\endhead
\bottomrule\noalign{}
\endlastfoot
Government & \(K_1\): restrictive/contraction policy & \(K_2\):
expansion policy & \(x\): share choosing \(K_1\) \\
Firms & \(M_1\): hire foreign talent & \(M_2\): hire domestic talent &
\(y\): share choosing \(M_1\) \\
Talent & \(N_1\): migrate to Japan & \(N_2\): alternative
destination/outside option & \(z\): share choosing \(N_1\) \\
\end{longtable}

Thus \(q=(x,y,z)\in[0,1]^3\). In particular, \textbf{\(x=0\) denotes
expansion}, while \(x=1\) denotes restriction. This orientation matters
when interpreting a trajectory toward \((0,1,1)\).

Time and payoffs were illustrative model units. The parameter values
were not estimated treatment effects, calibrated currency amounts or
observations of real migration rates.

\section{Model A: government-benefit
spillovers}\label{model-a-government-benefit-spillovers}

\subsection{Expected payoffs}\label{expected-payoffs}

The supplied expected payoffs were

\[
\begin{aligned}
\pi_G(K_1)&=E_2-C_2,\\
\pi_G(K_2)&=E_1-C_1-yzS_1-zS_2+R_1yz+R_2z,\\
\pi_F(M_1)&=-C_3+zE_3+(1-x)zS_1,\\
\pi_F(M_2)&=E_4-C_4,\\
\pi_T(N_1)&=E_5-C_5+(1-x)S_2,\\
\pi_T(N_2)&=E_6.
\end{aligned}
\]

\(S_1\) was a firm subsidy and \(S_2\) a talent subsidy under expansion.
In this specification the firm subsidy was conditional on foreign hiring
\textbf{and} migration, hence \(yzS_1\) in the government budget.
\(R_1yz\) and \(R_2z\) represented government benefits associated with
hiring/migration. They were distinct from a reduction in a firm's cost.

An equivalent pure-strategy payoff table follows. Entries are ordered as
(government, firm, talent).

\begin{longtable}[]{@{}
  >{\raggedright\arraybackslash}p{(\linewidth - 6\tabcolsep) * \real{0.2500}}
  >{\raggedright\arraybackslash}p{(\linewidth - 6\tabcolsep) * \real{0.2500}}
  >{\raggedright\arraybackslash}p{(\linewidth - 6\tabcolsep) * \real{0.2500}}
  >{\raggedright\arraybackslash}p{(\linewidth - 6\tabcolsep) * \real{0.2500}}@{}}
\toprule\noalign{}
\begin{minipage}[b]{\linewidth}\raggedright
Profile
\end{minipage} & \begin{minipage}[b]{\linewidth}\raggedright
Government
\end{minipage} & \begin{minipage}[b]{\linewidth}\raggedright
Firm
\end{minipage} & \begin{minipage}[b]{\linewidth}\raggedright
Talent
\end{minipage} \\
\midrule\noalign{}
\endhead
\bottomrule\noalign{}
\endlastfoot
\(K_1,M_1,N_1\) & \(E_2-C_2\) & \(E_3-C_3\) & \(E_5-C_5\) \\
\(K_1,M_1,N_2\) & \(E_2-C_2\) & \(-C_3\) & \(E_6\) \\
\(K_1,M_2,N_1\) & \(E_2-C_2\) & \(E_4-C_4\) & \(E_5-C_5\) \\
\(K_1,M_2,N_2\) & \(E_2-C_2\) & \(E_4-C_4\) & \(E_6\) \\
\(K_2,M_1,N_1\) & \(E_1-C_1-S_1-S_2+R_1+R_2\) & \(E_3-C_3+S_1\) &
\(E_5-C_5+S_2\) \\
\(K_2,M_1,N_2\) & \(E_1-C_1\) & \(-C_3\) & \(E_6\) \\
\(K_2,M_2,N_1\) & \(E_1-C_1-S_2+R_2\) & \(E_4-C_4\) & \(E_5-C_5+S_2\) \\
\(K_2,M_2,N_2\) & \(E_1-C_1\) & \(E_4-C_4\) & \(E_6\) \\
\end{longtable}

\subsection{Replicator equations}\label{replicator-equations}

For a two-strategy population with first-strategy share \(u\), the
average payoff is \(\bar\pi=u\pi_1+(1-u)\pi_2\). Therefore

\[\dot u=u(\pi_1-\bar\pi)=u(1-u)(\pi_1-\pi_2).\]

Define \(A=(E_2-C_2)-(E_1-C_1)\) and \(d=E_5-C_5-E_6\). The three
selection gaps are

\[
\begin{aligned}
G(y,z)&=A+(S_1-R_1)yz+(S_2-R_2)z,\\
F(x,z)&=-C_3-E_4+C_4+z[E_3+(1-x)S_1],\\
H(x)&=d+(1-x)S_2.
\end{aligned}
\]

The continuous-time system is

\[\dot x=x(1-x)G,\qquad \dot y=y(1-y)F,\qquad \dot z=z(1-z)H.\]

Each face of the unit cube is invariant in the exact mutation-free ODE:
when one share is exactly zero or one, its derivative is zero. A small
\emph{positive} share can nevertheless become extremely small before
recovering.

\section{Equilibria and local stability in Model
A}\label{equilibria-and-local-stability-in-model-a}

\subsection{Stationarity and the
Jacobian}\label{stationarity-and-the-jacobian}

The full stationarity condition is the simultaneous factorization

\[
\begin{gathered}
x\in\{0,1\}\ \text{or}\ G(y,z)=0,\\
y\in\{0,1\}\ \text{or}\ F(x,z)=0,\\
z\in\{0,1\}\ \text{or}\ H(x)=0.
\end{gathered}
\]

It includes eight corners and potentially mixed equilibria on edges,
faces or in the interior. A corner catalogue alone is not a general
enumeration of every equilibrium.

Writing \(a=x(1-x)\), \(b=y(1-y)\) and \(c=z(1-z)\), the analytic
Jacobian is

\[
J=\begin{pmatrix}
(1-2x)G & a(S_1-R_1)z & a[(S_1-R_1)y+S_2-R_2]\\
-bS_1z & (1-2y)F & b[E_3+(1-x)S_1]\\
-cS_2 & 0 & (1-2z)H
\end{pmatrix}.
\]

At a hyperbolic equilibrium, strictly negative eigenvalue real parts
imply local asymptotic stability; any positive real part implies
instability. A zero real part requires further analysis. These are local
dynamical statements, not a proof of global convergence or evolutionary
stability in every game-theoretic sense.

\subsection{The expansion--foreign-hiring--migration
corner}\label{the-expansionforeign-hiringmigration-corner}

At \(q^*=(0,1,1)\), the Jacobian is diagonal, with eigenvalues

\[
\begin{aligned}
\lambda_x&=A+S_1-R_1+S_2-R_2,\\
\lambda_y&=C_3+E_4-C_4-E_3-S_1,\\
\lambda_z&=-(d+S_2).
\end{aligned}
\]

Consequently, strict local stability requires all three conditions:

\[
\begin{aligned}
R_1+R_2&>A+S_1+S_2,\\
C_3&<E_3+S_1-(E_4-C_4),\\
S_2&>C_5+E_6-E_5.
\end{aligned}
\]

For the supplied baseline these give eigenvalues \((-0.7,-0.6,-0.1)\), a
hiring-cost bound \(C_3<1.5\), and a talent-subsidy bound \(S_2>0.5\).
These thresholds are conditional on this corner and specification. They
are not empirical policy estimates, and they must not be transferred to
Model B.

\subsection{A fully interior
candidate}\label{a-fully-interior-candidate}

When the denominators are nonzero, solving \(G=F=H=0\) gives

\[
\begin{aligned}
x^*&=1+\frac{d}{S_2},\\
z^*&=\frac{C_3+E_4-C_4}{E_3+(1-x^*)S_1},\\
y^*&=\frac{-A/z^*-(S_2-R_2)}{S_1-R_1}.
\end{aligned}
\]

All three coordinates must lie strictly between zero and one. The
baseline candidate has \(x^*=1/6\), \(z^*=0.825\), and
\(y^*\approx-0.803\), so it is infeasible. Degenerate denominators
require separate analysis.

At any fully interior stationary point of this particular mutation-free
specification, all three diagonal Jacobian entries vanish, so
\(\operatorname{tr}J=0\). Strict Hurwitz stability is therefore
impossible there. This observation does not decide nonhyperbolic
nonlinear stability.

\section{Model B: firm-cost spillovers and
mutation}\label{model-b-firm-cost-spillovers-and-mutation}

\subsection{Different payoff
assumptions}\label{different-payoff-assumptions}

The later source script used

\[
\begin{aligned}
C_3^{\mathrm{eff}}(z)&=\max(C_3-\eta z,0),\\
\pi_G(K_1)&=E_2-C_2,\\
\pi_G(K_2)&=E_1-C_1-yS_1-zS_2,\\
\pi_F(M_1)&=0.15E_4+zE_3+(1-x)S_1-C_3^{\mathrm{eff}}(z),\\
\pi_F(M_2)&=E_4-C_4,\\
\pi_T(N_1)&=E_5-C_5+(1-x)S_2,\qquad \pi_T(N_2)=E_6.
\end{aligned}
\]

Here the government paid \(yS_1\), rather than \(yzS_1\); the firm
received \((1-x)S_1\), rather than \((1-x)zS_1\); government benefits
\(R_1,R_2\) were absent; and the firm had an additional \(0.15E_4\)
term. That coefficient was preserved from the script, not estimated in
this release.

\subsection{Replicator dynamics with
mutation}\label{replicator-dynamics-with-mutation}

Let \(\Delta=(G_B,F_B,H)\) denote the payoff gaps derived from those
payoffs. For \(q_i\in\{x,y,z\}\),

\[\dot q_i=r_iq_i(1-q_i)\Delta_i+\mu_i(1/2-q_i).\]

The mutation term is linear mean reversion toward one half. At a zero
face it contributes \(+\mu_i/2\); at a one face it contributes
\(-\mu_i/2\). With positive mutation, the corners are \textbf{not}
equilibria and the exact vector field points inward on every face.

Away from the cost kink, write \(D=\partial\Delta/\partial q\). Then

\[
J_B=\operatorname{diag}(r_i(1-2q_i)\Delta_i-\mu_i)
+\operatorname{diag}(r_iq_i(1-q_i))D,
\]

where

\[
D=\begin{pmatrix}
0&S_1&S_2\\
-S_1&0&E_3+\eta\,\mathbf{1}_{\{C_3>\eta z\}}\\
-S_2&0&0
\end{pmatrix}.
\]

At \(C_3=\eta z\) the maximum is nondifferentiable. The implementation
refuses a smooth Jacobian there instead of silently choosing a
derivative. The supplied baseline and scan ranges remain on the
positive-effective-cost branch throughout the unit cube.

\clearpage
\section{Parameter dictionary and
baselines}\label{parameter-dictionary-and-baselines}

\begin{longtable}[]{@{}
  >{\raggedright\arraybackslash}p{(\linewidth - 6\tabcolsep) * \real{0.2143}}
  >{\raggedright\arraybackslash}p{(\linewidth - 6\tabcolsep) * \real{0.2143}}
  >{\raggedleft\arraybackslash}p{(\linewidth - 6\tabcolsep) * \real{0.2857}}
  >{\raggedleft\arraybackslash}p{(\linewidth - 6\tabcolsep) * \real{0.2857}}@{}}
\toprule\noalign{}
\begin{minipage}[b]{\linewidth}\raggedright
Parameter
\end{minipage} & \begin{minipage}[b]{\linewidth}\raggedright
Interpretation
\end{minipage} & \begin{minipage}[b]{\linewidth}\raggedleft
Model A
\end{minipage} & \begin{minipage}[b]{\linewidth}\raggedleft
Model B
\end{minipage} \\
\midrule\noalign{}
\endhead
\bottomrule\noalign{}
\endlastfoot
\(E_1,E_2\) & Expansion / restriction government benefits & 2.0, 1.5 &
2.60, 1.70 \\
\(C_1,C_2\) & Expansion / restriction government costs & 0.6, 0.2 &
0.35, 0.20 \\
\(E_3\) & Benefit conditional on foreign talent arriving & 2.0 & 1.20 \\
\(E_4,C_4\) & Domestic-hiring benefit and cost & 1.8, 0.5 & 1.15,
0.60 \\
\(C_3\) & Foreign-hiring friction cost & 0.9 & 1.05 \\
\(E_5,C_5\) & Japan migration benefit and cost & 2.0, 0.6 & 1.85,
0.80 \\
\(E_6\) & Outside-option payoff & 1.9 & 1.35 \\
\(S_1,S_2\) & Firm/talent subsidies & 0.8, 0.6 & 0.10, 0.22 \\
\(R_1,R_2\) & Government spillover benefits & 1.2, 0.8 & absent \\
\(\eta\) & Hiring-cost spillover coefficient & absent & 0.25 \\
\(r_x,r_y,r_z\) & Adaptation speeds & 1, 1, 1 & 1, 1, 1 \\
\(\mu_x,\mu_y,\mu_z\) & Mutation coefficients & 0, 0, 0 &
\((5,10,5)\times10^{-4}\) \\
\end{longtable}

Both source suites used \(T=200\), \(h=0.05\), and a final-30\% summary
window. Cases started at \((0.6,0.3,0.2)\), \((0.6,0.75,0.2)\) and
\((0.3,0.75,0.6)\). The policy case lowered \(C_3\) to 0.6 in Model A
and to 0.90 in Model B. Their different baselines and policy changes
must remain explicit.

\section{Numerical schemes and
validation}\label{numerical-schemes-and-validation}

The historical Model A solver used classical RK4:

\[q_{n+1}=\operatorname{clip}_{[\epsilon,1-\epsilon]}
\left[q_n+\frac h6(k_1+2k_2+2k_3+k_4)\right],\qquad \epsilon=10^{-9},\]

with clipped intermediate states inside the vector field. Model B used

\[q_{n+1}=\operatorname{clip}_{[\epsilon,1-\epsilon]}[q_n+h f_B(q_n)].\]

RK4 and Euler have different integration errors; the clipping operator
changes the numerical process further. A floor is not the same mechanism
as the explicitly specified mutation term.

An independent SciPy DOP853 reference was added for this release. For
Model A, positivity was protected using logit coordinates
\(u_i=\log(q_i/(1-q_i))\). Since \(\dot q_i=q_i(1-q_i)\Delta_i\), the
transformed equation is simply

\[\dot u_i=\Delta_i(\operatorname{logistic}(u)),\qquad q_i=\operatorname{logistic}(u_i).\]

This avoids replacing extremely small positive shares with an arbitrary
floor. A direct absolute-tolerance solver can otherwise lose their sign
and fail when the advantage later changes sign. Model B's positive
mutation kept its baseline equilibrium away from the boundary, so its
reference used the original state coordinates.

Verification comprised:

\begin{itemize}
\tightlist
\item
  matching original strategy payoffs and vector fields;
\item
  checking both analytic Jacobians with central finite differences,
  including both smooth cost branches;
\item
  checking Model A's eight corners and the eigenvalues at \((0,1,1)\);
\item
  checking inward mutation flux and invalid parameter/grid inputs;
\item
  comparing source-scheme paths with adaptive references;
\item
  reproducing all five archived tables: two case summaries, two \(C_3\)
  scans and one \(S_2\) scan;
\item
  checking the long-horizon clipping effect explicitly.
\end{itemize}

The frozen dependencies and machine-readable diagnostics are included in
\texttt{requirements.txt} and \texttt{results/validation.json}. Unit
tests do not replace calibration or a complete equilibrium analysis.

\section{Results and interpretation}\label{results-and-interpretation}

\subsection{Faithful reproduction of source
outputs}\label{faithful-reproduction-of-source-outputs}

The historical clipped government baseline ended near \((0,1,1)\) at
\(T=200\). The mutation model ended near
\((0.000334,0.000378,0.003112)\). Its small positive values reflected
mutation, rather than exact extinction. The English code reproduced the
archived output tables within their saved precision.

The source Model A hiring-cost grid contained 23 points from 0.3 to 1.4;
Model B contained 17 points from 0.6 to 1.4 and 13 subsidy points from 0
to 0.6. A plotted scan was a finite grid of simulations, not proof of a
unique bifurcation threshold.

The Model B baseline talent gap obeyed \(H(x)=-0.30+0.22(1-x)<0\)
throughout the cube. Increasing talent support or reducing
foreign-hiring costs therefore addressed different channels. The
baseline was not evidence that every subsidy necessarily produces high
adoption.

\subsection{Clipping versus the unmodified
ODE}\label{clipping-versus-the-unmodified-ode}

For the same Model A baseline and initial state, the source floor and
the continuous reference produced very different finite-horizon hiring
shares:

\begin{longtable}[]{@{}lr@{}}
\toprule\noalign{}
Method & \(y(200)\), approximately \\
\midrule\noalign{}
\endhead
\bottomrule\noalign{}
\endlastfoot
Historical RK4 with floor \(10^{-9}\) & 0.999999999 \\
Unmodified ODE, logit DOP853 & \(1.22\times10^{-21}\) \\
\end{longtable}

The continuous path reached a foreign-hiring share of approximately
\(1.54\times10^{-51}\) before recovering. Replacing that rare strategy
with a floor of \(10^{-9}\) dramatically shortened recovery. This
difference persisted when the historical time step was refined, because
reducing a time step does not remove the floor.

The locally stable \((0,1,1)\) corner and a long delay before reaching
it are compatible: local eigenvalues do not set a uniform arrival time
from every initial state. Neither finite-horizon output alone proved
global convergence. The repository therefore retained both the faithful
original numerical outputs and the continuous-model diagnostic, rather
than presenting the former as an accurate unclipped ODE solution.

The original case named ``path dependence'' changed the initial state.
That label alone did not establish multiple persistent attractors or
genuine asymptotic path dependence. Initial-state sensitivity, transient
delay and different attracting equilibria are distinct claims.

\section{Reproduction commands and accompanying
artifacts}\label{reproduction-commands-and-accompanying-artifacts}

From the repository root, using Python 3.12 or later:

\begin{Shaded}
\begin{Highlighting}[]
\ExtensionTok{python} \AttributeTok{{-}m}\NormalTok{ venv .venv}
\BuiltInTok{source}\NormalTok{ .venv/bin/activate}
\ExtensionTok{python} \AttributeTok{{-}m}\NormalTok{ pip install }\AttributeTok{{-}r}\NormalTok{ requirements.txt}
\ExtensionTok{python} \AttributeTok{{-}m}\NormalTok{ pytest }\AttributeTok{{-}q}
\ExtensionTok{python} \AttributeTok{{-}m}\NormalTok{ immigration\_game.reproduce}
\end{Highlighting}
\end{Shaded}

The public repository contains the English figures, CSV results, source
manifest and validation diagnostics:

\href{https://github.com/Kai1024-star/evolutionary-game-immigration-policy}{Evolutionary
Game Modeling of Immigration Policy}.

The figures separate government spillovers, firm-cost spillovers,
original-grid sensitivity and the clipping diagnostic. They are
generated from illustrative payoffs, not measured policy data. The
supplied MATLAB work and original manuscript, if available separately,
would require their own documentation and verification.

\end{document}
